\section{Materiais e Métodos}

O presente projeto será desenvolvido nas dependências do Laboratório de Otimização e Combinatória (LOCo) do Instituto de Computação, um ambiente compartilhado com outros estudantes de mestrado e doutorado que propicia a troca de ideias e a colaboração científica. Os algoritmos propostos serão testados empiricamente. Devido ao grande volume de instâncias disponíveis na literatura e aos elevados limites de tempo comumente adotados em trabalhos anteriores (frequentemente estabelecidos em uma hora por execução), as baterias de testes podem demandar dias para serem concluídas. Por essa razão, os experimentos computacionais serão executados nos servidores de alto desempenho disponibilizados pelo LOCo.

Para o suporte teórico e o levantamento do estado da arte, o candidato utilizará o acervo bibliográfico físico da biblioteca do Instituto de Matemática, Estatística e Computação Científica, além do Portal de Periódicos da CAPES, que assegura acesso institucional e gratuito às principais revistas científicas da área para os estudantes da Unicamp.

A metodologia de acompanhamento incluirá reuniões entre o candidato e o orientador para debater o andamento da pesquisa, superar desafios metodológicos e alinhar novas estratégias de resolução. Adicionalmente, encontros com os demais membros do grupo de pesquisa do orientador ocorrerão de forma periódica ou sob demanda, promovendo a integração e a discussão coletiva dos resultados alcançados.